\chapter{اللقاء الثامن: تابع تفسير سورة البقرة}

\section{تنبيهات منهجية في دراسة الكتب التأسيسية}
السلام عليكم ورحمة الله وبركاته. الحمد لله رب العالمين، وأشهد أن لا إله إلا الله وحده لا شريك له، وأشهد أن محمداً عبده ورسوله. 

أحب أن أنبه في البداية إلى أمر منهجي في دراستنا؛ نحن كنا نسير بالتركيز على كتاب واحد حتى ننهي منه جزءاً، وقد اقترح بعض الطلاب أن نجمع أنفسنا على كتاب واحد خلال الأسبوع (مثل أن ننهي سورة البقرة في الطبري) لتتهيأ لجو الكتاب وتعيش في مناخه، فهذا يسهل عليك التحصيل بدلاً من تشتيت الذهن بين عدة مصنفات.

كما أود التنبيه على أننا في دراستنا لتفسير \rawi{الطبري} سنختصر سرد الأسانيد الطويلة؛ لأن الغرض الأسمى هو الوقوف على المادة التفسيرية والنكت المنهجية، فسنقول اختصاراً: (روى الطبري بإسناده عن ابن عباس كذا...)، لتوفير الوقت للتعليق على ما قد يُشكل من الفهم.

\separator

\section{تأويل (في قلوبهم مرض)}
\isnad{قال الإمام الطبري رحمه الله:} «القول في تأويل قوله جل ثناؤه: \begin{ayah}فِي قُلُوبِهِم مَّرَضٌ\end{ayah}. واصل المرض السقم... وأخبر الله جل ثناؤه أن في قلوب المنافقين مرضاً، وإنما عنى جل ثناؤه بخبره عن مرض قلوبهم، الخبر عن مرض ما في قلوبهم من الاعتقاد. ولكن لما كان معلوماً بالخبر عن مرض القلب أنه معنيّ به مرض ما هم معتقدوه... استغنى بالخبر عن القلب بذلك».

ثم استشهد \rawi{الطبري} بأشعار العرب كقول \rawi{عمرو بن لجأ}: «وسبحت المدينة لا تلومها» يريد وسبح أهل المدينة.

\begin{fawaid}[تأثر الطبري بمدارس النحويين المعتزلة في تقدير المحذوف]
نلاحظ هنا فكرة ستتكرر معنا كثيراً؛ وهي تأثر \rawi{الطبري} رحمه الله ببعض النحويين (مثل: سيبويه، وأبي عبيدة معمر بن المثنى، والفراء، والأخفش) في مسألة افتراض قياس نحوي لكلام العرب، وجعل ما يخرج عنه محتاجاً للتقدير بالحذف أو الزيادة.

القرآن لم يقل: "في اعتقاد قلوبهم مرض"، بل قال: \begin{ayah}فِي قُلُوبِهِم مَّرَضٌ\end{ayah}. والصحابة لم يستشكلوا ذلك، بل فسروه بالشك والنفاق. ولكن اللغويين افترضوا أن القلب لا يمرض بهذا المعنى المادي، فقَدَّروا كلمة "اعتقاد". ومثل ذلك قولهم في \begin{ayah}وَاسْأَلِ الْقَرْيَةَ\end{ayah} أن الأصل: "اسأل أهل القرية". وهذا التوجه هو الذي فتح الباب لاحقاً للمتكلمين والمعتزلة لتحريف نصوص الصفات بدعوى المجاز، كقولهم في \begin{ayah}وَجَاءَ رَبُّكَ\end{ayah}: أي جاء "أمر" ربك. فيجب على طالب العلم أن يتفطن لمصادر هذه الأقوال ولا يسلم بها مباشرة.
\end{fawaid}

\separator

\section{تأويل (فزادهم الله مرضاً ولهم عذاب أليم بما كانوا يَكْذِبُونَ)}
\isnad{قال أبو جعفر:} «المرض الذي أخبر الله أنه زادهم على مرضهم، هو نظير ما كان في قلوبهم من الشك والحيرة... فزادهم الله بما أحدث من حدوده وفرائضه التي لم يكن فرضها... شكاً وحيرة».

\begin{fawaid}[أسباب الإضلال وتفسير القرآن بالقرآن]
هذه الآية تبين أن إضلال الله للعبد يكون بسببٍ من العبد، فالمنافقون كان في قلوبهم مرض (نفاق) فزادهم الله مرضاً فوقه. وهذا أحسن ما يُفسر به القرآن بالقرآن، كقوله تعالى: \begin{ayah}وَأَمَّا الَّذِينَ فِي قُلُوبِهِم مَّرَضٌ فَزَادَتْهُمْ رِجْسًا إِلَىٰ رِجْسِهِمْ وَمَاتُوا وَهُمْ كَافِرُونَ\end{ayah}.
\end{fawaid}

\subsection{الترجيح بين القراءات: (يَكْذِبُونَ) و (يُكَذِّبُونَ)}
\isnad{قال الطبري:} «اختلفت القراءة في ذلك؛ فقرأه بعضهم (يَكْذِبُونَ) مخففة... وقرأه آخرون (يُكَذِّبُونَ) بضم الياء وتشديد الذال... وكأن الذين قرأوا بالتشديد رأوا أن الله إنما أوجب العذاب لتكذيبهم نبيه... وليس الأمر في ذلك عندي كالذي قالوا؛ وذلك أن الله أنبأ عن المنافقين في أول النبأ عنهم بأنهم يكذبون بدعواهم الإيمان وأظهارهم ذلك بألسنتهم خداعاً... فأولى في حكمة الله أن يكون الوعيد منه لهم على ما افتتح به الخبر عنهم من قبيح أفعالهم (وهو الكذب)».

\begin{fawaid}[ترجيح القراءات المتواترة بالسياق]
لا حرج في الترجيح بين القراءات المتواترة لا لرد إحداها، بل لبيان أيهما أبلغ في إظهار المعنى المراد في سياق الآية. وقد رجح \rawi{الطبري} قراءة التخفيف (يَكْذِبُونَ)؛ لأن السياق من أول السورة يصف نفاقهم وادعاءهم الإيمان بألسنتهم وهم يبطنون الكفر، فالوعيد على هذا الكذب أليق بالسياق من الوعيد على مجرد التكذيب المطلق.
\end{fawaid}

\separator

\section{تأويل (وإذا قيل لهم لا تفسدوا في الأرض)}
\isnad{قال الطبري رحمه الله بعد إيراد الأقوال:} «وأولى التأويلين بالآية تأويل من قال إن قول الله... نزلت في المنافقين الذين كانوا على عهد رسول الله صلى الله عليه وسلم، وإن كان معنياً بها كل من كان بمثل صفتهم من المنافقين بعدهم إلى يوم القيامة».

\begin{fawaid}[الفرق بين نزول الآية وتنزّلها]
من الدقة في التفسير التفريق بين "النزول" (السياق التاريخي الذي نزلت فيه الآية، كنزولها في منافقي المدينة)، وبين "التنزل" (وهو سحب حكم الآية على كل من شابههم في الفعل عبر الأزمان).
\end{fawaid}

\subsection{مسألة العذر بالجهل في قوله (ولكن لا يشعرون)}
\isnad{قال الطبري:} «وذلك من حكم الله فيهم، أدلُ الدليل على تكذيب جَلَّ ثناؤه قول القائلين: إن عقوبات الله لا يستحقها إلا المعاند ربه... بعد علمه وثبوت الحجة عليه».

\begin{fawaid}[الموقف من العذر بالجهل ومناقشة الطبري]
كلام \rawi{الطبري} هنا قد يُفهم منه أنه ينفي العذر بالجهل، ويثبت العقوبة على من لا يشعر أنه على باطل. والصواب المحكم في الشريعة أن الله لا يُعذب أحداً حتى تبلغه الحجة ويعلم، كما قال تعالى: \begin{ayah}وَمَا كُنَّا مُعَذِّبِينَ حَتَّىٰ نَبْعَثَ رَسُولًا\end{ayah}.
أما معنى \begin{ayah}يَحْسَبُونَ أَنَّهُم مُّهْتَدُونَ\end{ayah} أو \begin{ayah}لَّا يَشْعُرُونَ\end{ayah}، فليس المراد أنهم كانوا يجهلون أمر النبي ولم تقم عليهم الحجة، بل الحجة قامت وهم يعرفون مخالفتهم، لكنهم "حسبوا" أن فعلهم هذا (وهو النفاق والتقية) هو الأذكى والأهدى لمصالحهم الدنيوية، فهم لا يشعرون بأنهم يُهلكون أنفسهم في الآخرة بهذا المكر.
\end{fawaid}

\separator

\section{تأويل (وإذا قيل لهم آمنوا كما آمن الناس)}
\isnad{قال الطبري في توجيهها:} «يعني وإذا قيل لهؤلاء الذين وصفهم الله... صدّقوا بمحمد صلى الله عليه وسلم وبما جاء به من عند الله كما صدّق به الناس... وإنما أدخلت الألف واللام في (الناس)... لأنهم كانوا معروفين عند الذين خوطبوا بذلك بأعيانهم (يعني الصحابة)».

\subsection{تأويل (قالوا أنؤمن كما آمن السفهاء)}
\isnad{قال الطبري:} «والسفيه: الجاهل الضعيف الرأي، القليل المعرفة بمواضع المنافع والمضار، ولذلك سمى الله جل وعز النساء والصبيان سفهاء، فقال: \begin{ayah}وَلَا تُؤْتُوا السُّفَهَاءَ أَمْوَالَكُمُ\end{ayah}... إنما عنى المنافقون بقيلهم: أنؤمن كما آمن السفهاء؟ إذ دعوا إلى التصديق بمحمد... قالوا: أنؤمن كما آمن أهل الجهل؟».

\begin{fawaid}[رد شبهة الطاعنين حول وصف النساء بالسفه]
يتصيد بعض أهل الريب هذا التفسير اللغوي لاتهام المفسرين بانتقاص المرأة. والمعنى الشرعي واللغوي أن "السفيه" هو من لا يُحسن التصرف في المال بغض النظر عن جنسه. وليس المعنى أن كل النساء فيهن هذه الصفة بالضرورة المطلقة، بل هي موعظة إرشادية لحفظ الأموال لمن عُرف بضعف تدبيره.
\end{fawaid}

\separator

\section{تأويل (وإذا لقوا الذين آمنوا قالوا آمنا)}
\isnad{قال الطبري:} «فإن قال قائل: أرأيت قوله \textit{وإذا خلوا إلى شياطينهم} فكيف قيل خلوا \textbf{إلى} شياطينهم؟ ولم يقل خلوا \textbf{بشياطينهم}... قيل: كان بعض نحويي البصرة يقول: يقال خلوت إلى فلان إذا أريد به خلوت إليه في الحاجة الخاصة... وبعض نحويي الكوفة كان يتأول أن ذلك بمعنى: وإذا انصرف خلاؤهم إلى شياطينهم».

\begin{fawaid}[تكلف اللغويين واختراع الإشكالات]
الصحابة وأئمة التابعين لم يستشكلوا حرف (إلى) هنا؛ لأن الكلام جارٍ على سعة لسان العرب. ولكن النحاة (كأبي عبيدة والفراء) لما افترضوا أن الأصل هو حرف (الباء)، جعلوا مجيء (إلى) إشكالاً يحتاج لتأويل بتضمين الفعل (أي انصرفوا) أو إنابة الحروف. وهذا التكلف هو الذي أطال الخلاف وجلب الشبهات للقرآن الكريم.
\end{fawaid}

\separator

\section{تأويل (الله يستهزئ بهم) والرد على المعطلة}
\isnad{قال أبو جعفر:} «اختلف في صفة استهزاء الله تعالى ذكره الذي ذكر أنه فاعله بالمنافقين... فقال بعضهم: الاستهزاء بهم توبيخه إياهم ولومه... أو إهلاكه إياهم... وقال آخرون: قوله (يخادعون الله وهو خادعهم) على الجواب (أي المجازاة)... وأن المكر والهزء حاق بهم».

وقد رد \rawi{الطبري} بقوة على المعتزلة الذين نفوا هذه الصفات عن الله بدعوى التنزيه، \isnad{قائلاً:} «وأما الذين زعموا أن قول الله (والله يستهزئ بهم) إنما هو على وجه الجواب، وأنه لم يك من الله استهزاء ولا مكر ولا خديعة، فنفوا عن الله جل ثناؤه ما قد أثبته لنفسه وأوجبه لها... يقال لهم: إن الله أخبرنا أنه مكر بقوم مضوا قبلنا... فما برهانكم على تفريقكم بزعمكم أنه خسف بمن أخبر أنه خسف به، ولم يمكر بمن أخبر أنه مكر به؟».

\begin{fawaid}[إثبات صفات الله على الكمال والرد على المعطلة]
هذا الموضع من أنفس المواضع في إثبات صفات الأفعال لله (كالمكر والاستهزاء والخديعة) على وجه الكمال والمجازاة بالعدل، دون تأويلها بالعبث واللعب. فكما أن الله ينتقم ويُعاقب حقيقةً، فهو يستهزئ بالمنافقين حقيقةً استهزاءً يليق بجلاله، ولا يصح نفي ما أثبته الله لنفسه فراراً من لوازم التشبيه بالمخلوقين.
\end{fawaid}

\separator

\section{تأويل (يمدهم في طغيانهم يعمهون)}
\isnad{قال أبو جعفر:} «وأولى هذه الأقوال بالصواب في قوله (ويمدهم) أن يكون بمعنى: يزيدهم على وجه الإملاء والترك لهم في عتوهم وتمردهم... كما قال تعالى: \begin{ayah}وَنَذَرُهُمْ فِي طُغْيَانِهِمْ يَعْمَهُونَ\end{ayah}. والطغيان: الفعلان من طغى إذا تجاوز في الأمر حده. ويعمهون: من العمه وهو الضلال، يترددون حيارى».

\section{تأويل (أولئك الذين اشتروا الضلالة بالهدى)}
\isnad{قال أبو جعفر:} «والذي هو أولى عندي بتأويل الآية... أخذوا الضلالة وتركوا الهدى. وذلك أن كل كافر بالله فإنه مستبدل بالإيمان كفراً، وذلك هو معنى الشراء؛ لأن كل مشتري شيئاً فإنما يستبدل مكان الذي يؤخذ منه بدلاً آخر... فما ربحت تجارتهم: أي خسروا ولم يربحوا باختيارهم الحيرة على الرشاد».

\separator

\section{تأويل (مثلهم كمثل الذي استوقد ناراً)}
\isnad{طرح الطبري إشكالاً نحوياً ثم أجاب عنه:} «كيف قيل: مثلهم كمثل \textbf{الذي} استوقد ناراً؟ وقد علمت أن الهاء والميم من (مثلهم) كناية جماع، و(الذي) دلالة على واحد؟ ... قيل: إنما جاز لأن المراد من الخبر عن مكر المنافقين، الخبر عن مثل استضاءتهم بما أظهروا من الإقرار... والاستضاءة وإن اختلفت أشخاص أهلها معنى واحد».

ثم بين أن الله ضرب المثل للمنافقين الذين أظهروا الإيمان بألسنتهم فعصموا دماءهم وأموالهم (وهذا هو النور المؤقت)، ولكنهم لما أبطنوا الكفر، أذهب الله نورهم في الآخرة وتركهم في ظلمات النفاق يعمهون.

\section{تأويل (أو كصيب من السماء)}
\isnad{قال أبو جعفر:} «والصيب: الفيعل من قولك: صاب المطر يصوب صوباً، إذا انحدر ونزل... وتأويل ذلك: مثل استضاءة المنافقين بضوء إقرارهم بالإسلام مع استسرارهم الكفر، كمثل غيث سرى ليلاً في مزنة ظلماء، يحدوها رعد ويستطير في حافاتها برق شديد اللمعان».

وهذا مثل ثانٍ، ضربه الله للمنافقين؛ فالبرق اللامع يمثل ظاهر إيمانهم الذي يمشون في ضوئه (ينتفعون به دنيوياً)، والرعد والصواعق تمثل وعيد الله والزواجر التي في القرآن والتي تُفزعهم وتكشف نفاقهم، فيجعلون أصابعهم في آذانهم حذر الفضيحة وحذر الموت.

\subsection{توجيه (ولو شاء الله لذهب بسمعهم وأبصارهم)}
\isnad{قال أبو جعفر:} «وإنما خص الله السمع والأبصار بأنه لو شاء أذهبها من المنافقين دون سائر جوارحهم؛ للذي جرى من ذكرها في الآيتين... وعيداً من الله لهم، ووصفاً لنفسه أنه المقتدر عليهم، ليتقوا بأسه ويسارعوا إليه بالتوبة، إن الله على كل شيء قدير».

وبهذا ينتهي اللقاء الثامن بحمد الله، ونكمل في اللقاء القادم بإذن الله.